\section{From \texorpdfstring{$\cxor$}{cxor} back to statistical compression}\label{sec:ladder}

$\cxor_{\W5}$ can be read as a context model reduced to one bit, and its limitations are exactly
what that reduction removes.  A cell can say ``zero'' or ``one,'' but not ``one with probability
$0.9$.''  It cannot distinguish a prediction supported by one observation from one supported by
a thousand.  $\W5$ changes its mind after every error, has no escape when a context is new, and
cannot share evidence between related contexts.

Putting capacity back gives a useful conceptual ladder:
\[
\begin{gathered}
\boxed{\text{one-bit $\cxor$ predictor}}\\[-1pt]
\big\downarrow\\[-1pt]
\boxed{\text{successor counts and probabilities}}\\[-1pt]
\big\downarrow\\[-1pt]
\boxed{\text{PPM with escape/backoff}}\\[-1pt]
\big\downarrow\\[-1pt]
\boxed{\text{context mixing or learned prediction}}
\end{gathered}
\]
The first step replaces a hard remembered successor by evidence about both possible successors.
PPM combines context orders through an explicit escape mechanism \cite{bell1990text}.  Context
mixing combines predictions from many models, often with adaptive weights
\cite{Mahoney2005AdaptiveWO}.  These are not variants among the sixteen wirings: widening the
state leaves the finite one-bit design space.

What survives is the causal architecture.  At position $k$, a deterministic predictor whose
state is reproducible from $x_0,\dots,x_{k-1}$ may produce a hard decision $\hat x_k$, and
\[
  r_k=x_k+\hat x_k
\]
remains an invertible residual for exactly the reason used in
Theorem~\ref{thm:causal-inverse}.  The predictor may contain counts, mixtures, or learned state;
the decoder recovers $x_k=r_k+\hat x_k$ and then performs the same state update.  Exact
reproducibility, including numerical arithmetic, is essential.

Probabilistic compressors normally do something more informative.  They pass the probability
$P(x_k=1\mid x_{<k})$ to an arithmetic or range coder rather than threshold it to one bit and
discard confidence.  $\cxor_{\W5}$ isolates the degenerate hard-prediction case, in which correctness alone
is represented.  It is useful analytically because its state and complete update space are small,
not because one bit is expected to compete with a statistical model.

The practical question left by the example in Section~\ref{sec:intro} is correspondingly narrow:
for which sources, context orders, and downstream coders does contextual differencing expose
structure that the back end models more cheaply than the original sequence?  Answering it needs
comparative experiments that include the model description, initialisation, coder, and any side
information.  Zero enrichment alone is a diagnostic of prediction success, not a compression
result.
